\section{Spectral gap}\label{sec:spectral_gap}

A frustration-free Hamiltonian is \emph{gapped} if the first excited
eigenvalue is bounded away from zero uniformly in the chain length $N$.
For MPS parent Hamiltonians, the spectral gap follows from the martingale
method of Fannes--Nachtergaele--Werner and
Nachtergaele~\cite{Fannes1992Finitely, Nachtergaele1996Spectral}, as recalled
by Cirac--Perez-Garcia--Schuch--Verstraete~\cite[Section~IV.C]{Cirac2021Matrix}
and in the later formulation of Kastoryano and
Lucia~\cite{Kastoryano2018Martingale}.

\begin{definition}[Parent ground space in the Hilbert-space realization]
\label{def:parent_ground_space_es}
    \lean{MPSTensor.parentHamiltonianGroundSpaceES}
    \leanok
    \uses{def:parent_hamiltonian, def:nsite_space_parent}
    Denote by $\mathcal G_{N,L}^{\mathrm{ES}}(A)$ the parent ground space after
    identifying the coefficient space with the Hilbert space
    $\ell^{2}(\{0,\ldots,d-1\}^N)$.
\end{definition}

\begin{definition}[Parent Hamiltonian in the Hilbert-space realization]
\label{def:parent_hamiltonian_es}
    \lean{MPSTensor.parentHamiltonianES}
    \leanok
    \uses{def:parent_hamiltonian, def:nsite_space_parent}
    Denote by $H_{N,L}^{\mathrm{ES}}(A)$ the conjugate of $H_{N}(A,L)$ by the
    canonical identification between the coefficient space and
    $\ell^{2}(\{0,\ldots,d-1\}^N)$.
\end{definition}

\begin{theorem}[Kernel of the Hilbert-space parent Hamiltonian]
    \label{thm:parent_ground_space_es_kernel}
    \lean{MPSTensor.parentHamiltonianGroundSpaceES_eq_ker_parentHamiltonianES}
    \leanok
    Under the canonical $\ell^2$ identification of the $N$-site state space, the
    Hilbert-space parent ground space is exactly the kernel of the corresponding
    parent Hamiltonian.
\end{theorem}

\begin{proof}
    \leanok
    This is an immediate unraveling of the definitions: both constructions are
    obtained from the same parent Hamiltonian by conjugating with the canonical
    identification between the site space and the Hilbert-space realization, so
    the ground space is precisely the kernel after this identification.
\end{proof}

\begin{lemma}[Membership in the Hilbert-space parent ground space]
    \label{lem:mem_parent_ground_space_es}
    \lean{MPSTensor.mem_parentHamiltonianGroundSpaceES_iff}
    \leanok
    \uses{thm:parent_ground_space_es_kernel}
    A vector belongs to the Hilbert-space parent ground space exactly when the
    parent Hamiltonian annihilates it.
\end{lemma}

\begin{remark}[Euclidean local projector ingredients]
    \label{rem:euclidean_local_projector_ingredients}
    \lean{MPSTensor.parentInteractionES}
    \lean{MPSTensor.cyclicRestrictES}
    \lean{MPSTensor.localTermES}
    \lean{MPSTensor.localTermESSummand}
    \leanok
    \uses{def:ground_space_es, rem:cyclic_window_operators}
    Let $P_L^{\mathrm{ES}}(A)$ denote the orthogonal projector onto
    $(G_L^{\mathrm{ES}}(A))^\perp$ on the Euclidean $L$-site space. For a cyclic
    window starting at $i$ and an outside configuration $\tau$, the transported
    restriction map $R_{i,\tau}$ evaluates an $N$-site Euclidean vector on that
    window. The transported local term $h_{i,\mathrm{ES}}$ is the corresponding
    $N$-site local projector, and $R_{i,\tau}^{\dagger} P_L^{\mathrm{ES}}(A) R_{i,\tau}$ is
    the basic positive summand in its cyclic-averaging formula.
\end{remark}

\begin{lemma}[Positivity of the Euclidean local summands]
    \label{lem:euclidean_local_projector_positive}
    \lean{MPSTensor.parentInteractionES_isPositive}
    \lean{MPSTensor.localTermESSummand_isPositive}
    \lean{MPSTensor.localTermESSummand_sum_isPositive}
    \leanok
    \uses{rem:euclidean_local_projector_ingredients}
    The Euclidean local projector $P_L^{\mathrm{ES}}(A)$ is positive. Consequently,
    every summand $R_{i,\tau}^{\dagger} P_L^{\mathrm{ES}}(A) R_{i,\tau}$ is positive, and so
    is the finite sum over $\tau$.
\end{lemma}

\begin{proof}
    \leanok
    Orthogonal projectors are positive. Conjugation by $R_{i,\tau}$ preserves
    positivity, and finite sums of positive operators remain positive.
\end{proof}

\begin{lemma}[Cyclic averaging for the transported ES local terms]
    \label{lem:local_term_es_average}
    \lean{MPSTensor.parentHamiltonianES_eq_sum_localTermES}
    \lean{MPSTensor.localTermES_eq_average_localTermESSummand}
    \leanok
    \uses{rem:euclidean_local_projector_ingredients, lem:euclidean_local_projector_positive}
    The transported parent Hamiltonian is the sum of the transported local terms
    $H_{N,\mathrm{ES}} = \sum_i h_{i,\mathrm{ES}}$, and each transported local term admits
    the exact cyclic-averaging formula
    \begin{equation}\label{eq:ph_cyclic_averaging}
        h_{i,\mathrm{ES}}
        = d^{-L} \sum_{\tau} R_{i,\tau}^{\dagger} P_L^{\mathrm{ES}}(A) R_{i,\tau}.
    \end{equation}
\end{lemma}

\begin{proof}
    \leanok
    Unfold the transported local term and compare it configuration-by-configuration
    with the cyclic restrictions. The $d^{-L}$ factor compensates for the $d^L$
    choices of the window coordinates in the ambient configuration parameter $\tau$.
\end{proof}

\begin{lemma}[Positivity of the transported ES Hamiltonian]
    \label{lem:transported_es_positive}
    \lean{MPSTensor.localTermES_isPositive}
    \lean{MPSTensor.parentHamiltonianES_isPositive}
    \leanok
    \uses{lem:local_term_es_average}
    Each transported local term $h_{i,\mathrm{ES}}$ is positive, hence the full
    transported parent Hamiltonian $H_{N,\mathrm{ES}}$ is positive.
\end{lemma}

\begin{proof}
    \leanok
    Apply the averaging formula~(\ref{eq:ph_cyclic_averaging}): $h_{i,\mathrm{ES}}$ is a non-negative scalar
    multiple of a finite sum of positive conjugates of $P_L^{\mathrm{ES}}(A)$, hence is
    itself positive. Summing over $i$ yields positivity of $H_{N,\mathrm{ES}}$.
\end{proof}

\begin{lemma}[Symmetric-projection structure of transported ES local terms]
    \label{lem:transported_es_local_projections}
    \lean{MPSTensor.parentInteractionES_isSymmetricProjection}
    \lean{MPSTensor.localTermES_isSymmetricProjection}
    \leanok
    \uses{rem:euclidean_local_projector_ingredients, lem:transported_es_positive}
    The Euclidean local projector $P_L^{\mathrm{ES}}(A)$ is a symmetric projection, and
    every transported local term $h_{i,\mathrm{ES}}$ on the $N$-site chain is a
    symmetric projection.
\end{lemma}

\begin{proof}
    \leanok
    \uses{lem:transported_es_positive}
    The local operator $P_L^{\mathrm{ES}}(A)$ is an orthogonal projection onto
    $(G_L^{\mathrm{ES}}(A))^\perp$. For $L\le N$, restricting
    $h_{i,\mathrm{ES}}v$ to a cyclic window gives
    $P_L^{\mathrm{ES}}(A)$ applied to the corresponding restriction of $v$; applying
    $h_{i,\mathrm{ES}}$ a second time therefore applies
    $(P_L^{\mathrm{ES}}(A))^2=P_L^{\mathrm{ES}}(A)$ on that window. The positivity
    result above supplies symmetry, and the case $L>N$ is the zero projection by
    definition.
\end{proof}

\begin{lemma}[Kernel of the Euclidean parent interaction]\label{lem:parent_interaction_es_kernel}
    \lean{MPSTensor.parentInteractionES_apply_eq_zero_iff}
    \leanok
    \uses{def:ground_space_es, rem:euclidean_local_projector_ingredients}
    The Euclidean local interaction $P_L^{\mathrm{ES}}(A)$ annihilates exactly the
    Euclidean local ground space:
    \[
        P_L^{\mathrm{ES}}(A)v=0 \quad\Longleftrightarrow\quad
        v\in G_L^{\mathrm{ES}}(A).
    \]
\end{lemma}

\begin{proof}
    \leanok
    Since $P_L^{\mathrm{ES}}(A)$ is the orthogonal projector onto
    $(G_L^{\mathrm{ES}}(A))^\perp$, its kernel is the original subspace
    $G_L^{\mathrm{ES}}(A)$.
\end{proof}

\begin{lemma}[Local-term kernel via cyclic restrictions]\label{lem:local_term_es_kernel_restrictions}
    \lean{MPSTensor.localTermES_eq_zero_iff_forall_cyclicRestrictES_mem_groundSpaceES}
    \lean{MPSTensor.cyclicRestrictES_mem_groundSpaceES_of_localTermES_eq_zero}
    \leanok
    \uses{lem:parent_interaction_es_kernel, rem:euclidean_local_projector_ingredients}
    For $L\le N$, a transported local term $h_{i,\mathrm{ES}}$ annihilates a vector
    $v$ iff every boundary-filled restriction of $v$ to the cyclic $L$-window
    starting at $i$ lies in $G_L^{\mathrm{ES}}(A)$.
\end{lemma}

\begin{proof}
    \leanok
    Restricting $h_{i,\mathrm{ES}}v$ to the same cyclic window gives
    $P_L^{\mathrm{ES}}(A)$ applied to the corresponding restriction of $v$.
    Lemma~\ref{lem:parent_interaction_es_kernel} identifies the vanishing of this
    local projector with membership in $G_L^{\mathrm{ES}}(A)$.
\end{proof}

\begin{theorem}[Adjacent local kernels realize the intersection property]
    \label{thm:adjacent_local_kernels_ground_space_es}
    \lean{MPSTensor.mem_groundSpaceES_succ_of_adjacent_localTermES_eq_zero}
    \lean{MPSTensor.adjacent_localTermES_eq_zero_of_mem_groundSpaceES_succ}
    \lean{MPSTensor.adjacent_localTermES_eq_zero_iff_mem_groundSpaceES_succ}
    \leanok
    \uses{thm:ground_space_intersection, lem:local_term_es_kernel_restrictions,
        lem:mem_ground_space_es}
    If $A$ is injective and $L\ge2$, then on an $(L+1)$-site chain the kernels of
    the two adjacent transported $L$-site local terms are exactly the transported
    $(L+1)$-site MPS ground space:
    \[
        h_{0,\mathrm{ES}}v=0\ \text{ and }\ h_{1,\mathrm{ES}}v=0
        \quad\Longleftrightarrow\quad
        v\in G_{L+1}^{\mathrm{ES}}(A).
    \]
\end{theorem}

\begin{proof}
    \leanok
    \uses{thm:ground_space_intersection, lem:local_term_es_kernel_restrictions,
        lem:mem_ground_space_es}
    Lemma~\ref{lem:local_term_es_kernel_restrictions} translates the two local
    kernel assumptions into the left and right $L$-site ground conditions. The
    open-chain intersection property, Theorem~\ref{thm:ground_space_intersection},
    then grows back the shared $(L-1)$-site overlap to the $(L+1)$-site ground space.
    The converse uses the forward restriction lemmas for ground-space vectors and
    the same kernel--restriction characterization.
\end{proof}

\begin{lemma}[Non-overlap positivity for transported ES local terms]
    \label{lem:transported_es_nonoverlap_positive}
    \lean{MPSTensor.CyclicWindowsDisjoint}
    \lean{MPSTensor.CyclicWindowsDisjoint.symm}
    \lean{MPSTensor.CyclicWindowsDisjoint.of_not_cyclicWindowsOverlap}
    \lean{MPSTensor.localTermES_commute_of_cyclic_windows_disjoint}
    \lean{MPSTensor.localTermES_re_inner_nonneg_of_cyclic_windows_disjoint}
    \leanok
    \uses{lem:transported_es_local_projections, rem:projection_geometry_rowsum_identities}
    Let $L\le N$ and let $h_{i,\mathrm{ES}},h_{j,\mathrm{ES}}$ be transported local
    terms whose cyclic $L$-windows are site-disjoint: no site of the periodic chain
    has cyclic offset $<L$ from both $i$ and $j$. Then the terms commute pointwise:
    for every vector $v$,
    \[
        h_{i,\mathrm{ES}}(h_{j,\mathrm{ES}}v)
        = h_{j,\mathrm{ES}}(h_{i,\mathrm{ES}}v),
    \]
    and their ordered cross term is non-negative:
    \[
        0\le \operatorname{Re}\langle h_{i,\mathrm{ES}}v,
        h_{j,\mathrm{ES}}v\rangle.
    \]
\end{lemma}

\begin{proof}
    \leanok
    \uses{lem:transported_es_local_projections, rem:projection_geometry_rowsum_identities}
    In coordinates, $h_{i,\mathrm{ES}}$ applies $P_L^{\mathrm{ES}}(A)$ only to the
    $i$-window variables and leaves the complementary variables fixed. Site
    disjointness gives two replacement identities: replacing the $i$-window does
    not change the extracted $j$-window, and the $i$- and $j$-window replacements
    commute. Therefore the two embedded local operators commute. The preceding
    projection-geometry identity for commuting symmetric projections gives the
    non-negative ordered cross term.
\end{proof}

\begin{lemma}[Quadratic form to norm bound]\label{lem:quadratic_form_to_norm_bound}
    \lean{FrustrationFree.spectralGap_of_martingale}
    \leanok
    Let $H$ be a positive semidefinite self-adjoint operator on a
    finite-dimensional complex Hilbert space, and suppose that, for all~$v$,
    \[
        \gamma\,\langle v, H v\rangle \le \langle H v, H v\rangle
    \]
    i.e.\ $H^2 \ge \gamma H$ as a quadratic form, for some $\gamma > 0$.
    Then for every $v$ in the orthogonal complement of $\ker H$,
    \[
        \gamma\,\|v\| \le \|H v\|.
    \]
    In particular, every positive eigenvalue of $H$ is at least $\gamma$.
    The positive-semidefiniteness hypothesis is essential: without it,
    $H = -\mathrm{Id}$ with $\gamma = 2$ satisfies the quadratic-form
    inequality vacuously but violates the norm bound.
\end{lemma}

\begin{proof}
    \leanok
    By the finite-dimensional spectral theorem, $H$ diagonalises in an
    orthonormal eigenbasis with real eigenvalues $\lambda_k$. For a unit
    eigenvector $\psi_k$ with eigenvalue $\lambda_k$ the hypothesis
    gives $\gamma\,\lambda_k \le \lambda_k^2$, whence $\lambda_k \le 0$
    or $\lambda_k \ge \gamma$. Since $H$ is positive semidefinite
    (its quadratic form is $\ge 0$), the negative case is excluded and
    every nonzero eigenvalue satisfies $\lambda_k \ge \gamma$. Writing
    $v = \sum_k c_k \psi_k$ and projecting onto $(\ker H)^{\perp}$
    keeps only terms with $\lambda_k \ge \gamma$, so
    $\|H v\|^2 = \sum_k |c_k|^2 \lambda_k^2
        \ge \gamma^2 \sum_k |c_k|^2 = \gamma^2 \|v\|^2$.
\end{proof}

\begin{lemma}[Parent-Hamiltonian quadratic-form reduction]
    \label{lem:parent_hamiltonian_es_quadratic_reduction}
    \lean{MPSTensor.parentHamiltonianES_norm_bound_of_quadratic_form}
    \lean{MPSTensor.parentHamiltonianES_gap_bound_of_quadratic_form}
    \leanok
    \uses{lem:transported_es_positive, thm:parent_ground_space_es_kernel,
        lem:quadratic_form_to_norm_bound}
    Let $H_{N,L}^{\mathrm{ES}}(A)$ be the transported parent Hamiltonian. If a
    constant $\gamma > 0$ satisfies the uniform quadratic-form estimate
    \begin{equation}\label{eq:ph_uniform_quadratic_form}
        \gamma\,\operatorname{Re}\langle H_{N,L}^{\mathrm{ES}}(A)v, v\rangle
        \le
        \operatorname{Re}\langle H_{N,L}^{\mathrm{ES}}(A)v,
        H_{N,L}^{\mathrm{ES}}(A)v\rangle
    \end{equation}
    for every $N \ge 2L$ and every vector $v$, then
    \[
        \gamma\,\|v\| \le \|H_{N,L}^{\mathrm{ES}}(A)v\|
    \]
    for every $v \in (\mathcal G_{N,L}^{\mathrm{ES}}(A))^\perp$. In particular,
    with $L>1$ and $\gamma = 1/(4L)$, the explicit gap-bound conclusion follows
    once this same quadratic-form estimate~(\ref{eq:ph_uniform_quadratic_form}) is supplied.
\end{lemma}

\begin{proof}
    \leanok
    The transported Hamiltonian is positive by
    Lemma~\ref{lem:transported_es_positive}, and its transported ground space is
    its kernel by Theorem~\ref{thm:parent_ground_space_es_kernel}. Therefore the
    abstract spectral-theorem conversion of
    Lemma~\ref{lem:quadratic_form_to_norm_bound} applies directly. The
    explicit constant is positive because $L>1$ implies $4L>0$.
\end{proof}

\begin{remark}[Projection-geometry identities for row-sum arguments]
    \label{rem:projection_geometry_rowsum_identities}
    \lean{LinearMap.IsSymmetricProjection.re_inner_apply_apply_self}
    \lean{LinearMap.IsSymmetricProjection.re_inner_nonneg}
    \lean{LinearMap.IsSymmetricProjection.re_inner_apply_apply_nonneg_of_commute}
    \lean{LinearMap.IsSymmetricProjection.re_inner_apply_apply_ge_neg_of_norm_apply_le}
    \lean{ProjectionGeometry.re_inner_sum_apply_left}
    \lean{ProjectionGeometry.re_inner_sum_apply_apply}
    \lean{ProjectionGeometry.sum_sum_eq_diag_add_offdiag}
    \lean{ProjectionGeometry.crossTerm_sum_bound_of_ordered_rowSum}
    \leanok
    The projection-geometry reduction uses the following elementary Hilbert-space
    identities: the diagonal identity
    $\operatorname{Re}\langle Pv,Pv\rangle=\operatorname{Re}\langle Pv,v\rangle$
    for symmetric projections; nonnegativity of this quadratic form; nonnegativity
    of $\operatorname{Re}\langle Pv,Qv\rangle$ for commuting symmetric projections
    $P,Q$; the Cauchy--Schwarz conversion from
    $\|P Q v\|\le c\|Pv\|$ to
    $\operatorname{Re}\langle Pv,Qv\rangle\ge -c\operatorname{Re}\langle Pv,v\rangle$;
    finite-sum expansions of
    $\operatorname{Re}\langle (\sum_i P_i)v,v\rangle$ and
    $\operatorname{Re}\langle (\sum_i P_i)v,(\sum_i P_i)v\rangle$; the
    diagonal/off-diagonal split of the ordered double sum; and the aggregate
    off-diagonal estimate obtained from row sums $\sum_{j\ne i}c_{ij}\le1$.
\end{remark}

\begin{lemma}[Ordered projection row-sum quadratic-form reduction]
    \label{lem:ordered_projection_rowsum_reduction}
    \lean{ProjectionGeometry.quadraticForm_sum_projections_of_ordered_rowSum}
    \leanok
    \uses{rem:projection_geometry_rowsum_identities}
    Let $P_i$ be a finite family of symmetric projections and let
    $H=\sum_i P_i$. If the ordered off-diagonal forms obey
    \[
        \operatorname{Re}\langle P_i v, P_j v\rangle
        \ge -(1-\gamma)c_{ij}\operatorname{Re}\langle P_i v,v\rangle
        \qquad (i\ne j),
    \]
    with row sums $\sum_{j\ne i} c_{ij}\le 1$ and $\gamma\le 1$, then
    $H^2\ge \gamma H$ as a quadratic form.
\end{lemma}

\begin{proof}
    \leanok
    Expand $\operatorname{Re}\langle H v,Hv\rangle$ as an ordered double sum.
    Symmetric idempotence identifies each diagonal term with
    $\operatorname{Re}\langle P_i v,v\rangle$, while the row-summable ordered
    cross-term bounds control the off-diagonal contribution by
    $-(1-\gamma)\sum_i\operatorname{Re}\langle P_i v,v\rangle$.
\end{proof}

\begin{lemma}[Finite-overlap projection row-sum reduction]
    \label{lem:finite_overlap_projection_rowsum_reduction}
    \lean{ProjectionGeometry.quadraticForm_sum_projections_of_finite_overlap}
    \leanok
    \uses{lem:ordered_projection_rowsum_reduction}
    Let $P_i$ be a finite family of symmetric projections and let $\gamma\le1$.
    Suppose an interaction predicate $i\sim j$ has at most $m$ off-diagonal
    neighbours in each row for some positive integer $m$. If noninteracting
    ordered cross terms are non-negative and, on each interacting pair,
    \[
        \operatorname{Re}\langle P_i v,P_jv\rangle
        \ge -(1-\gamma)\frac{1}{m}\operatorname{Re}\langle P_i v,v\rangle,
    \]
    then $H=\sum_i P_i$ satisfies $H^2\ge \gamma H$ as a quadratic form.
\end{lemma}

\begin{proof}
    \leanok
    Choose $c_{ij}=1/m$ on interacting pairs and $c_{ij}=0$ otherwise. The row
    cardinality bound gives $\sum_{j\ne i}c_{ij}\le1$; noninteracting
    nonnegativity supplies the zero-coefficient cross-term estimates. Apply
    Lemma~\ref{lem:ordered_projection_rowsum_reduction}.
\end{proof}

\begin{lemma}[Finite-overlap projection norm-compression reduction]
    \label{lem:finite_overlap_projection_norm_reduction}
    \lean{ProjectionGeometry.quadraticForm_sum_projections_of_finite_overlap_norm_bound}
    \lean{ProjectionGeometry.quadraticForm_sum_projections_of_finite_overlap_norm_bound_of_le}
    \leanok
    \uses{lem:finite_overlap_projection_rowsum_reduction,
        rem:projection_geometry_rowsum_identities}
    Let $P_i$ be a finite family of symmetric projections and let $\gamma\le1$.
    Suppose an interaction predicate has at most $m>0$ off-diagonal neighbours
    in each row, noninteracting ordered cross terms are non-negative, and
    interacting pairs satisfy the norm-compression estimate
    \begin{equation}\label{eq:ph_finite_overlap_norm_compression}
        \|P_iP_jv\|\le (1-\gamma)\frac{1}{m}\|P_iv\|.
    \end{equation}
    Then $H=\sum_iP_i$ satisfies $H^2\ge\gamma H$ as a quadratic form. The same
    conclusion holds from a separate compression coefficient $\eta$ whenever
    $\eta\le (1-\gamma)/m$.
\end{lemma}

\begin{proof}
    \leanok
    \uses{lem:finite_overlap_projection_rowsum_reduction,
        rem:projection_geometry_rowsum_identities}
    Cauchy--Schwarz and symmetric idempotence turn the norm-compression estimate
    into
    $\operatorname{Re}\langle P_i v,P_jv\rangle
        \ge -(1-\gamma)m^{-1}\operatorname{Re}\langle P_iv,v\rangle$.
    If a separate coefficient $\eta$ is used, the assumption
    $\eta\le(1-\gamma)/m$ first weakens that estimate to
    (\ref{eq:ph_finite_overlap_norm_compression}). The finite-overlap row-sum
    reduction then applies.
\end{proof}

\begin{lemma}[Parent-Hamiltonian ordered row-sum quadratic-form reduction]
    \label{lem:parent_hamiltonian_ordered_rowsum_quadratic}
    \lean{MPSTensor.parentHamiltonianES_quadratic_form_of_ordered_local_term_bounds}
    \leanok
    \uses{lem:ordered_projection_rowsum_reduction, def:parent_hamiltonian_es,
        rem:euclidean_local_projector_ingredients, lem:transported_es_local_projections}
    For a fixed chain length $N$, suppose the transported local terms
    $h_{i,\mathrm{ES}}$ satisfy the ordered row-sum estimate
    \begin{equation}\label{eq:ph_ordered_rowsum_estimate}
        \operatorname{Re}\langle h_{i,\mathrm{ES}} v,h_{j,\mathrm{ES}}v\rangle
        \ge -(1-\gamma)c_{ij}\operatorname{Re}\langle h_{i,\mathrm{ES}}v,v\rangle
        \qquad (i\ne j),
    \end{equation}
    with $\sum_{j\ne i}c_{ij}\le1$ and $\gamma\le1$. Then the transported
    parent Hamiltonian satisfies
    \[
        \gamma\operatorname{Re}\langle H_{N,L}^{\mathrm{ES}}v,v\rangle
        \le
        \operatorname{Re}\langle H_{N,L}^{\mathrm{ES}}v,
        H_{N,L}^{\mathrm{ES}}v\rangle
    \]
    for every vector $v$.
\end{lemma}

\begin{proof}
    \leanok
    \uses{lem:transported_es_local_projections, lem:ordered_projection_rowsum_reduction}
    Lemma~\ref{lem:transported_es_local_projections} supplies the symmetric-projection
    hypothesis for the family $P_i=h_{i,\mathrm{ES}}$. Apply
    Lemma~\ref{lem:ordered_projection_rowsum_reduction} to this family and use
    $H_{N,L}^{\mathrm{ES}}=\sum_i h_{i,\mathrm{ES}}$.
\end{proof}

\begin{lemma}[Parent-Hamiltonian gap from ordered local row bounds]
    \label{lem:parent_hamiltonian_ordered_rowsum_gap}
    \lean{MPSTensor.parentHamiltonianES_gap_bound_of_ordered_local_term_bounds}
    \leanok
    \uses{lem:parent_hamiltonian_ordered_rowsum_quadratic,
        lem:parent_hamiltonian_es_quadratic_reduction, lem:transported_es_local_projections}
    Assume uniformly for every $N\ge2L$ that the transported local terms satisfy
    the ordered row-sum estimate~(\ref{eq:ph_ordered_rowsum_estimate}) with $\gamma=1/(4L)$. If $L>1$, then
    $1/(4L)>0$ and, for every
    $v\in(\mathcal G_{N,L}^{\mathrm{ES}}(A))^\perp$,
    \[
        \frac{1}{4L}\|v\|\le
        \|H_{N,L}^{\mathrm{ES}}(A)v\|.
    \]
\end{lemma}

\begin{proof}
    \leanok
    \uses{lem:parent_hamiltonian_ordered_rowsum_quadratic,
        lem:parent_hamiltonian_es_quadratic_reduction}
    The fixed-chain row-sum reduction supplies the quadratic-form estimate with
    $\gamma=1/(4L)$ for every admissible $N$. The positivity and kernel
    identification of the transported Hamiltonian then allow
    Lemma~\ref{lem:parent_hamiltonian_es_quadratic_reduction} to convert this
    quadratic-form estimate into the norm lower bound on the orthogonal
    complement of the ground space.
\end{proof}

\begin{lemma}[Parent-Hamiltonian finite-overlap Friedrichs quadratic reduction]
    \label{lem:parent_hamiltonian_finite_overlap_friedrichs_quadratic}
    \lean{MPSTensor.parentHamiltonianES_quadratic_form_of_finite_overlap_friedrichs}
    \leanok
    \uses{lem:finite_overlap_projection_rowsum_reduction, def:parent_hamiltonian_es}
    For a fixed chain length $N$, let $m$ be a positive integer and let
    $\gamma\le1$. Suppose an overlap predicate on local windows has at most $m$
    overlapping off-diagonal windows in each row. If non-overlapping local terms
    have non-negative ordered cross terms, overlapping local terms satisfy the
    finite-overlap Friedrichs estimate
    \[
        \operatorname{Re}\langle h_{i,\mathrm{ES}}v,h_{j,\mathrm{ES}}v\rangle
        \ge -(1-\gamma)\frac{1}{m}
        \operatorname{Re}\langle h_{i,\mathrm{ES}}v,v\rangle,
    \]
    and the transported local terms are symmetric projections, then the
    transported Hamiltonian satisfies $H_{N,L}^{\mathrm{ES}}{}^2\ge
        \gamma H_{N,L}^{\mathrm{ES}}$ as a quadratic form.
\end{lemma}

\begin{proof}
    \leanok
    Apply the finite-overlap projection reduction with
    $P_i=h_{i,\mathrm{ES}}$.
\end{proof}

\begin{lemma}[Finite-overlap norm-compression reduction]
    \label{lem:parent_hamiltonian_finite_overlap_norm_quadratic}
    \lean{MPSTensor.parentHamiltonianES_quadratic_form_of_finite_overlap_norm_bound}
    \lean{MPSTensor.parentHamiltonianES_quadratic_form_of_finite_overlap_norm_bound_of_le}
    \leanok
    \uses{lem:finite_overlap_projection_norm_reduction, def:parent_hamiltonian_es,
        lem:transported_es_local_projections}
    For a fixed chain length $N$, let $m>0$ and $\gamma\le1$. Suppose an overlap
    predicate on local windows has at most $m$ overlapping off-diagonal windows in
    each row, non-overlapping local terms have non-negative ordered cross terms,
    and overlapping local terms satisfy
    \[
        \|h_{i,\mathrm{ES}}(h_{j,\mathrm{ES}}v)\|
        \le (1-\gamma)\frac{1}{m}\|h_{i,\mathrm{ES}}v\|.
    \]
    Then $H_{N,L}^{\mathrm{ES}}{}^2\ge \gamma H_{N,L}^{\mathrm{ES}}$ as a
    quadratic form. The same conclusion holds from a separate coefficient $\eta$
    satisfying $\eta\le(1-\gamma)/m$.
\end{lemma}

\begin{proof}
    \leanok
    \uses{lem:finite_overlap_projection_norm_reduction,
        lem:transported_es_local_projections}
    Apply the finite-overlap norm-compression projection reduction to the family
    $P_i=h_{i,\mathrm{ES}}$. If a separate coefficient $\eta$ is used, first weaken
    the compression estimate using $\eta\le(1-\gamma)/m$.
\end{proof}

\begin{definition}[Cyclic-window support]
    \label{def:cyclic_window_support}
    \lean{MPSTensor.cyclicWindowSupport}
    \leanok
    \uses{rem:cyclic_window_operators}
    For a length-$L$ cyclic window on $\Fin{N}$ starting at $i$, its support is
    \[
        S_i=\{i,i+1,\ldots,i+L-1\}\pmod N,
    \]
    with repeated visits counted once.
\end{definition}

\begin{lemma}[Cyclic-window support offsets]
    \label{lem:cyclic_window_support_offsets}
    \lean{MPSTensor.mem_cyclicWindowSupport_iff}
    \leanok
    \uses{def:cyclic_window_support}
    If $L\le N$, then a site $k$ belongs to the cyclic support $S_i$ exactly when
    its cyclic offset from $i$ is below $L$:
    \[
        k\in S_i
        \quad\Longleftrightarrow\quad
        (k-i)\bmod N < L.
    \]
\end{lemma}

\begin{proof}
    \leanok
    Write $k=i+r\pmod N$ for a representative offset $r$. Membership in
    $S_i$ gives an offset $r<L$, and conversely an offset below $L$ gives the
    corresponding point of the support.
\end{proof}

\begin{definition}[Cyclic-window overlap]
    \label{def:cyclic_windows_overlap}
    \lean{MPSTensor.cyclicWindowsOverlap}
    \leanok
    \uses{def:cyclic_window_support}
    Two length-$L$ cyclic windows overlap, written $i\sim_L j$, when their
    supports meet: $S_i\cap S_j\ne\varnothing$.
\end{definition}

\begin{lemma}[Concrete non-overlap positivity for cyclic supports]
    \label{lem:cyclic_window_nonoverlap_positive}
    \lean{MPSTensor.CyclicWindowsDisjoint.of_not_cyclicWindowsOverlap}
    \lean{MPSTensor.localTermES_re_inner_nonneg_of_not_cyclicWindowsOverlap}
    \leanok
    \uses{lem:cyclic_window_support_offsets, def:cyclic_windows_overlap,
        lem:transported_es_nonoverlap_positive}
    If $L\le N$ and the cyclic supports $S_i$ and $S_j$ do not meet, then the
    transported local ES terms have non-negative ordered cross term:
    \[
        0\le\operatorname{Re}\langle h_{i,\mathrm{ES}}v,
        h_{j,\mathrm{ES}}v\rangle.
    \]
\end{lemma}

\begin{proof}
    \leanok
    \uses{lem:cyclic_window_support_offsets, lem:transported_es_nonoverlap_positive}
    The support-offset characterization converts failure of $S_i\cap S_j\ne\varnothing$
    into site-disjointness of the two cyclic windows. The disjoint-window
    positivity lemma then applies.
\end{proof}

\begin{lemma}[Cyclic-window overlap row cardinality]
    \label{lem:cyclic_window_overlap_cardinality}
    \lean{MPSTensor.cyclicWindowsOverlap_card_le}
    \leanok
    \uses{def:cyclic_window_support, def:cyclic_windows_overlap}
    If $2L\le N$ and $L>1$, then for every $i\in\Fin{N}$,
    \[
        \#\{j\ne i: i\sim_L j\}\le 2(L-1).
    \]
\end{lemma}

\begin{proof}
    \leanok
    \uses{def:cyclic_window_support, def:cyclic_windows_overlap}
    If $S_i$ and $S_j$ meet, choose offsets $a,b<L$ with
    $i+a\equiv j+b\pmod N$. Since $2L\le N$, the offset from $i$ to $j$ is
    either the clockwise distance $a-b\in\{1,\ldots,L-1\}$ or the
    counterclockwise distance $b-a\in\{1,\ldots,L-1\}$; the zero distance is
    excluded by $j\ne i$. Thus every overlapping off-diagonal start lies among
    the $L-1$ clockwise starts or the $L-1$ counterclockwise starts.
\end{proof}

\begin{lemma}[Parent-Hamiltonian gap from finite-overlap Friedrichs data]
    \label{lem:parent_hamiltonian_finite_overlap_friedrichs_gap}
    \lean{MPSTensor.parentHamiltonianES_gap_bound_of_finite_overlap_friedrichs}
    \leanok
    \uses{lem:parent_hamiltonian_finite_overlap_friedrichs_quadratic,
        lem:parent_hamiltonian_es_quadratic_reduction}
    Fix $L>1$ and set $m=2(L-1)$. Suppose uniformly for every $N\ge2L$ that
    the hypotheses of the fixed-chain finite-overlap Friedrichs reduction hold
    with $\gamma=1/(4L)$ and this value of $m$. Then $1/(4L)>0$, and the
    explicit norm lower bound with constant $1/(4L)$ follows for vectors in
    $(\mathcal G_{N,L}^{\mathrm{ES}}(A))^\perp$.
\end{lemma}

\begin{proof}
    \leanok
    The fixed-chain finite-overlap reduction gives the required quadratic-form
    estimate for every admissible $N$. Lemma~\ref{lem:parent_hamiltonian_es_quadratic_reduction}
    converts that estimate to the norm lower bound.
\end{proof}

\begin{lemma}[Parent-Hamiltonian gap from cyclic-window Friedrichs data]
    \label{lem:parent_hamiltonian_cyclic_window_friedrichs_gap}
    \lean{MPSTensor.parentHamiltonianES_gap_bound_of_cyclic_window_friedrichs}
    \leanok
    \uses{lem:parent_hamiltonian_finite_overlap_friedrichs_gap,
        lem:cyclic_window_overlap_cardinality, lem:transported_es_local_projections,
        lem:cyclic_window_nonoverlap_positive}
    Fix $L>1$ and let $i\sim_L j$ mean that the cyclic supports of the
    length-$L$ windows at $i$ and $j$ meet. Suppose uniformly for every
    $N\ge2L$ that overlapping off-diagonal terms satisfy
    \begin{equation}\label{eq:ph_cyclic_window_overlap_estimate}
        \operatorname{Re}\langle h_{i,\mathrm{ES}}v,h_{j,\mathrm{ES}}v\rangle
        \ge -(1-\tfrac{1}{4L})\frac{1}{2(L-1)}
        \operatorname{Re}\langle h_{i,\mathrm{ES}}v,v\rangle.
    \end{equation}
    Then $1/(4L)>0$, and the explicit norm lower bound with constant $1/(4L)$
    follows for vectors in $(\mathcal G_{N,L}^{\mathrm{ES}}(A))^\perp$.
\end{lemma}

\begin{proof}
    \leanok
    \uses{lem:parent_hamiltonian_finite_overlap_friedrichs_gap,
        lem:cyclic_window_overlap_cardinality, lem:transported_es_local_projections,
        lem:cyclic_window_nonoverlap_positive}
    Lemma~\ref{lem:cyclic_window_overlap_cardinality} gives the row-cardinality
    bound $m=2(L-1)$ for the overlap relation among length-$L$ local windows,
    and
    Lemma~\ref{lem:transported_es_local_projections} supplies the symmetric
    projection structure of each transported local term. If two cyclic supports
    do not meet, the corresponding windows are site-disjoint, so
    Lemma~\ref{lem:cyclic_window_nonoverlap_positive} gives the required
    non-negative ordered cross term. Apply
    Lemma~\ref{lem:parent_hamiltonian_finite_overlap_friedrichs_gap} with these
    facts and the overlapping-window estimate~(\ref{eq:ph_cyclic_window_overlap_estimate}).
\end{proof}

\begin{lemma}[Parent-Hamiltonian gap from overlapping cyclic Friedrichs data]
    \label{lem:parent_hamiltonian_cyclic_overlap_friedrichs_gap}
    \lean{MPSTensor.parentHamiltonianES_gap_bound_of_cyclic_window_overlap_friedrichs}
    \leanok
    \uses{lem:parent_hamiltonian_cyclic_window_friedrichs_gap}
    Fix $L>1$ and let $i\sim_L j$ mean that the cyclic supports of the
    length-$L$ windows at $i$ and $j$ meet. Suppose uniformly for every
    $N\ge2L$ that overlapping off-diagonal terms satisfy
    \begin{equation}\label{eq:ph_cyclic_overlap_friedrichs}
        \operatorname{Re}\langle h_{i,\mathrm{ES}}v,h_{j,\mathrm{ES}}v\rangle
        \ge -(1-\tfrac{1}{4L})\frac{1}{2(L-1)}
        \operatorname{Re}\langle h_{i,\mathrm{ES}}v,v\rangle.
    \end{equation}
    Then the explicit norm lower bound with constant $1/(4L)$ follows for
    vectors in $(\mathcal G_{N,L}^{\mathrm{ES}}(A))^\perp$.
\end{lemma}

\begin{proof}
    \leanok
    \uses{lem:parent_hamiltonian_cyclic_window_friedrichs_gap}
    Apply Lemma~\ref{lem:parent_hamiltonian_cyclic_window_friedrichs_gap}.
    Its non-overlap positivity and finite-row hypotheses are already internal to
    the preceding reduction; the overlapping estimate
    (\ref{eq:ph_cyclic_overlap_friedrichs}) supplies its remaining ordered
    cross-term condition.
\end{proof}

\begin{lemma}[Parent-Hamiltonian gap from overlapping cyclic norm-compression data]
    \label{lem:parent_hamiltonian_cyclic_overlap_norm_gap}
    \lean{MPSTensor.parentHamiltonianES_gap_bound_of_cyclic_window_overlap_norm_bound}
    \leanok
    \uses{lem:parent_hamiltonian_cyclic_overlap_friedrichs_gap,
        rem:projection_geometry_rowsum_identities}
    Fix $L>1$. Suppose uniformly for every $N\ge2L$ and every overlapping
    off-diagonal pair that
    \[
        \|h_{i,\mathrm{ES}}(h_{j,\mathrm{ES}}v)\|
        \le (1-\tfrac{1}{4L})\frac{1}{2(L-1)}
        \|h_{i,\mathrm{ES}}v\|.
    \]
    Then the explicit norm lower bound with constant $1/(4L)$ follows for
    vectors in $(\mathcal G_{N,L}^{\mathrm{ES}}(A))^\perp$.
\end{lemma}

\begin{proof}
    \leanok
    \uses{lem:parent_hamiltonian_cyclic_overlap_friedrichs_gap,
        rem:projection_geometry_rowsum_identities}
    The norm-compression condition is converted to the ordered Friedrichs lower
    bound by the projection-geometry Cauchy--Schwarz lemma. The preceding
    overlapping-cyclic Friedrichs reduction then gives the norm lower bound.
\end{proof}

\begin{lemma}[Parent-Hamiltonian gap from a cyclic overlap compression constant]
    \label{lem:parent_hamiltonian_cyclic_overlap_norm_constant_gap}
    \lean{MPSTensor.parentHamiltonianES_gap_bound_of_cyclic_window_overlap_norm_bound_of_le}
    \lean{MPSTensor.parentHamiltonianES_gap_bound_of_cyclic_window_overlap_norm_bound_of_lt}
    \leanok
    \uses{lem:parent_hamiltonian_finite_overlap_norm_quadratic,
        lem:cyclic_window_overlap_cardinality, lem:cyclic_window_nonoverlap_positive,
        lem:parent_hamiltonian_es_quadratic_reduction}
    Fix $L>1$, and let $m=2(L-1)$. Suppose uniformly for every $N\ge2L$ and every
    overlapping off-diagonal pair that
    \[
        \|h_{i,\mathrm{ES}}(h_{j,\mathrm{ES}}v)\|
        \le \eta\,\|h_{i,\mathrm{ES}}v\|.
    \]
    If $\gamma>0$, $\gamma\le1$, and $\eta\le(1-\gamma)/m$, then $0<\gamma$
    and, for every admissible $N$ and every
    $v\in(\mathcal G_{N,L}^{\mathrm{ES}}(A))^\perp$,
    \[
        \gamma\|v\|\le \|H_{N,L}^{\mathrm{ES}}(A)v\|.
    \]
    The assertion $0<\gamma$ repeats the hypothesis, in the same shape as the
    strict statement below. In particular, if
    $\eta\ge0$ and $\eta m<1$, then $0<1-\eta m$ and, for every admissible $N$
    and every $v\in(\mathcal G_{N,L}^{\mathrm{ES}}(A))^\perp$,
    \[
        (1-\eta m)\|v\|\le \|H_{N,L}^{\mathrm{ES}}(A)v\|.
    \]
\end{lemma}

\begin{proof}
    \leanok
    \uses{lem:parent_hamiltonian_finite_overlap_norm_quadratic,
        lem:cyclic_window_overlap_cardinality, lem:cyclic_window_nonoverlap_positive,
        lem:parent_hamiltonian_es_quadratic_reduction}
    The overlap-cardinality lemma gives at most $m=2(L-1)$ overlapping
    off-diagonal neighbours in each row, while the non-overlap positivity lemma
    handles disjoint supports. The finite-overlap norm-compression reduction gives
    the quadratic-form estimate with the chosen $\gamma$, and the positive
    quadratic-form criterion converts it to the norm lower bound. The strict form
    is the same statement with $\gamma=1-\eta m$.
\end{proof}

\begin{theorem}[Martingale criterion for spectral gap]\label{thm:martingale_criterion}
    \notready
    \uses{def:is_frustration_free, lem:quadratic_form_to_norm_bound}
    Consider a frustration-free Hamiltonian
    \[
        H = \sum_{i=1}^n h_i
    \]
    with projector interaction terms ($h_i^2 = h_i$).
    Suppose that for every unordered pair $\{i,j\}$ of overlapping terms there
    exist constants $c_{\{i,j\}} \ge 0$ and $\gamma > 0$ satisfying
    \begin{equation}\label{eq:ph_martingale_condition}
        h_i h_j + h_j h_i
        \ge -c_{\{i,j\}}(1-\gamma)(h_i + h_j),
    \end{equation}
    and suppose further that the coefficients attached to the terms overlapping
    $h_i$ satisfy
    \[
        \sum_{\substack{j\ne i\\ \{i,j\}\text{ overlapping}}}c_{\{i,j\}}\le1
    \]
    for every $i$.
    Then the smallest positive eigenvalue of $H$ satisfies
    $\lambda_{\min}^+(H) \ge \gamma$.
\end{theorem}

\begin{proof}
    \notready
    The relevant geometry is a family of translated length-$L$ projector
    windows on the ring; two terms $h_i$ and $h_j$ interact only when their
    supports overlap, equivalently when the cyclic distance $d_N(i,j)$ is
    smaller than $L$.
    From $h_i^2 = h_i$,
    \[
        H^2 = \sum_i h_i + \sum_{\substack{i < j \\ \text{overlapping}}}
        (h_i h_j + h_j h_i) + \sum_{\substack{i < j \\ \text{non-overlapping}}}
        (h_i h_j + h_j h_i).
    \]
    Non-overlapping projectors commute ($h_i h_j = h_j h_i$), so their
    product $h_i h_j$ is itself a projector and hence PSD; thus the last
    sum is $\ge 0$.
    For the overlapping sum, the martingale condition~(\ref{eq:ph_martingale_condition}) gives
    $h_i h_j + h_j h_i \ge -c_{\{i,j\}}(1-\gamma)(h_i + h_j)$.
    Summing over unordered pairs $\{i,j\}$:
    \[
        \sum_{\{i,j\}, \text{overlap}} (h_i h_j + h_j h_i)
        \ge -(1-\gamma) \sum_{\{i,j\}, \text{overlap}}
        c_{\{i,j\}}(h_i + h_j).
    \]
    Collecting terms for each $h_i$:
    \[
        \sum_{\{i,j\}, \text{overlap}} c_{\{i,j\}}(h_i + h_j)
        = \sum_i \Bigl(\sum_{\substack{j \ne i \\ \text{overlap}}}
        c_{\{i,j\}}\Bigr) h_i
        \le H,
    \]
    where the last inequality uses the row-sum hypothesis
    $\sum_{\substack{j\ne i\\ \{i,j\}\text{ overlapping}}}c_{\{i,j\}}\le1$.
    Combining, $H^2 \ge H - (1-\gamma)H = \gamma H$.
    For any eigenvector with $H\psi = \lambda\psi$ and $\lambda > 0$,
    this gives $\lambda^2 \ge \gamma\lambda$, hence
    $\lambda \ge \gamma$~\cite{Kastoryano2018Martingale}.
\end{proof}

\begin{lemma}[Martingale condition for MPS]\label{lem:martingale_mps}
    \notready
    \uses{thm:martingale_criterion, def:parent_interaction,
        def:local_term_parent, thm:ground_space_intersection, def:injective}
    Let $A$ be injective and fix an interaction range $L \ge 2$ as in
    Definition~\ref{def:parent_interaction}. Let $h_L(A) := \Pi_{G_L(A)^\perp}$
    be the $L$-site parent interaction, and for each chain length and site $i$
    let $h_i$ denote its translate as in Definition~\ref{def:local_term_parent}.
    Then for every pair of overlapping terms $h_i$ and $h_j$ with
    $0 < d_N(i,j) < L$, where $d_N(i,j) := \min(|i-j|,\, N-|i-j|)$ is the
    cyclic distance on the ring (i.e.\ whose $L$-site supports share at
    least one site under periodic boundary conditions),
    the martingale condition of Theorem~\ref{thm:martingale_criterion} holds with
    constants $c_{\{i,j\}}$ and $\gamma > 0$ depending only on $A$ (not on $N$).
    Cirac--Perez-Garcia--Schuch--Verstraete~\cite[Section~IV.C]{Cirac2021Matrix}
    obtain this uniformity by the martingale method, invoking finite-chain-state
    estimates to control the constants. For the explicit cyclic-window
    constants displayed here, the remaining quantitative step is the
    overlapping-window estimate used as a hypothesis in
    Theorem~\ref{thm:friedrichs_gap_bound}; the preceding lemmas reduce the
    martingale inequality and the spectral gap to that estimate.
    (When $L = 1$ the interaction terms have disjoint supports, so there are no
    overlapping pairs and the spectral gap follows directly without the
    martingale argument.)
\end{lemma}

\begin{proof}
    \uses{thm:ground_space_intersection, def:injective}
    \notready
    Two terms $h_i, h_j$ with $0 < d_N(i,j) < L$ have supports overlapping
    in $L - d_N(i,j)$ sites (where $d_N$ is the cyclic distance on
    the $N$-site ring). Since $A$ is injective, the intersection property
    (Theorem~\ref{thm:ground_space_intersection}) gives the qualitative identity
    $\ker(h_i) \cap \ker(h_j) = G_{[i,j+L-1]}(A)$.
    Because the local Hilbert space is finite-dimensional, the subspaces
    $\ker(h_i)$ and $\ker(h_j)$ live in a fixed finite-dimensional space.
    The intersection identity shows that $\ker(h_i) \cap \ker(h_j)$ is
    exactly $G_{[i,j+L-1]}(A)$; in particular there is no non-trivial
    vector in one kernel that is orthogonal to the other yet lies in both.
    Hence the Friedrichs angle $\theta_{d_N(i,j)}$ between $\ker(h_i)$ and
    $\ker(h_j)$ is strictly positive (for the fixed injective tensor
    $A$). Because the angle depends only on the overlap pattern
    $d_N(i,j) \in \{1,\ldots,L-1\}$ and not on the absolute position,
    the number of distinct angles is at most $L-1$.
    Let $\alpha := \max_{1 \le s < L} \cos\theta_s < 1$ be the
    worst-case cosine of the Friedrichs angles.

    \emph{Row-sum bound.}
    Since $L \ge 2$, each site $i$ overlaps with at most $2(L-1)$
    neighbours (those $j$ with $0 < d_N(i,j) < L$).
    Set $c_{\{i,j\}} := \tfrac{1}{2(L-1)}$ for every unordered overlapping
    pair and $c_{\{i,j\}} := 0$ otherwise. These constants depend only on the
    cyclic distance $d_N(i,j)$.
    Then
    $\sum_{\substack{j\ne i\\ \{i,j\}\text{ overlapping}}}c_{\{i,j\}}
    \le 2(L-1)\cdot\tfrac{1}{2(L-1)} = 1$,
    satisfying the row-sum hypothesis of
    Theorem~\ref{thm:martingale_criterion}.

    \emph{Explicit cyclic-window constant.}
    With the row coefficient above and
    \[
        \gamma_L=\frac{1}{4L},
    \]
    the martingale condition is the fixed cyclic-window estimate
    \[
        h_i h_j + h_j h_i
        \ge
        -\left(1-\frac{1}{4L}\right)\frac{1}{2(L-1)}(h_i+h_j)
    \]
    for every overlapping off-diagonal pair. A sufficient form, used below, is
    the one-sided norm-compression estimate
    \[
        \|h_i h_j v\|
        \le
        \left(1-\frac{1}{4L}\right)\frac{1}{2(L-1)}\|h_i v\|.
    \]
    The qualitative Friedrichs-angle conclusion $\alpha<1$ does not by itself
    imply this numerical bound; the missing step is the uniform
    overlapping-window estimate with the displayed coefficient. Thus the
    remaining comparison is whether the principal-angle estimates cited in
    \cite{Fannes1992Finitely, Nachtergaele1996Spectral,
    Kastoryano2018Martingale} supply this cyclic-window inequality, with
    constants independent of $N$, under the convention used here.
    % Paper-gap note: docs/paper-gaps/cpgsv21_martingale_overlap.tex.
\end{proof}

\begin{theorem}[Gap bound from cyclic overlap compression]
    \label{thm:friedrichs_gap_bound}
    \lean{MPSTensor.parentHamiltonianES_gap_bound_of_friedrichs}
    \leanok
    \uses{def:parent_ground_space_es, def:parent_hamiltonian_es,
        thm:parent_ground_space_es_kernel,
        rem:euclidean_local_projector_ingredients,
        lem:euclidean_local_projector_positive,
        lem:local_term_es_average,
        lem:transported_es_local_projections,
        lem:transported_es_nonoverlap_positive,
        lem:parent_hamiltonian_es_quadratic_reduction,
        lem:parent_hamiltonian_ordered_rowsum_gap,
        lem:parent_hamiltonian_finite_overlap_friedrichs_quadratic,
        lem:parent_hamiltonian_finite_overlap_friedrichs_gap,
        lem:cyclic_window_overlap_cardinality,
        lem:cyclic_window_nonoverlap_positive,
        lem:parent_hamiltonian_cyclic_window_friedrichs_gap,
        lem:parent_hamiltonian_cyclic_overlap_friedrichs_gap,
        lem:parent_hamiltonian_cyclic_overlap_norm_gap}
    Let $A$ be injective and let $L>1$. Suppose that, for every $N\ge2L$ and
    every overlapping off-diagonal pair of cyclic length-$L$ windows,
    \[
        \|h_{i,\mathrm{ES}}(h_{j,\mathrm{ES}}v)\|
        \le
        \left(1-\frac{1}{4L}\right)\frac{1}{2(L-1)}
        \|h_{i,\mathrm{ES}}v\|.
    \]
    Then the parent Hamiltonian has the explicit gap constant
    \[
        \gamma_{L} := \frac{1}{4L} > 0.
    \]
    Concretely, if $N \ge 2L$ and
    $v \in (\mathcal G_{N,L}^{\mathrm{ES}}(A))^{\perp}$, then
    \[
        \gamma_{L} \|v\| \le \|H_{N,L}^{\mathrm{ES}}(A) v\|.
    \]
\end{theorem}

\begin{proof}
    \leanok
    The local averaging, positivity, and symmetric-projection statements are now
    established in Lemmas~\ref{lem:local_term_es_average},
    \ref{lem:transported_es_positive}, and~\ref{lem:transported_es_local_projections}.
    Lemma~\ref{lem:parent_hamiltonian_es_quadratic_reduction} reduces the present
    theorem to the remaining MPS-specific task: deriving the uniform
    quadratic-form estimate. Lemma~\ref{lem:parent_hamiltonian_ordered_rowsum_gap}
    establishes the finite-sum algebra from ordered cross-term row bounds, and
    Lemmas~\ref{lem:parent_hamiltonian_finite_overlap_friedrichs_quadratic}
    and~\ref{lem:parent_hamiltonian_finite_overlap_friedrichs_gap} state the
    common finite-overlap specialization with $m=2(L-1)$.
    Lemma~\ref{lem:transported_es_nonoverlap_positive} supplies the non-overlap
    positivity statement for site-disjoint windows, and
    Lemma~\ref{lem:cyclic_window_nonoverlap_positive} translates this to the
    concrete support-disjoint cyclic windows.
    Lemma~\ref{lem:cyclic_window_overlap_cardinality} gives the finite-overlap
    row-cardinality bound.
    Lemma~\ref{lem:parent_hamiltonian_cyclic_window_friedrichs_gap} combines
    these ingredients with the finite-overlap reduction, so the only additional
    input is the Friedrichs-angle estimate for overlapping cyclic windows.
    Lemma~\ref{lem:parent_hamiltonian_cyclic_overlap_friedrichs_gap} states the
    same overlap-only formulation, and
    Lemma~\ref{lem:parent_hamiltonian_cyclic_overlap_norm_gap} states a sufficient
    norm-compression formulation. Applying that norm-compression formulation to
    the displayed hypothesis gives the stated lower bound with
    $\gamma_L=1/(4L)$.
\end{proof}

\begin{theorem}[Conditional spectral gap for MPS parent Hamiltonians]
    \label{thm:parent_hamiltonian_gapped}
    \lean{MPSTensor.parentHamiltonian_gapped}
    \leanok
    \uses{thm:friedrichs_gap_bound}
    For an injective tensor $A$ and a range $L > 1$, assume the overlapping
    cyclic-window norm-compression estimate of
    Theorem~\ref{thm:friedrichs_gap_bound}. Then there exists $\gamma > 0$ such
    that for every $N \ge 2L$ and every vector
    $v \in (\mathcal G_{N,L}^{\mathrm{ES}}(A))^{\perp}$ one has
    \[
        \gamma \|v\| \le \|H_{N,L}^{\mathrm{ES}}(A) v\|.
    \]
    By Lemma~\ref{lem:quadratic_form_to_norm_bound}, this implies the usual
    uniform lower bound on every nonzero eigenvalue of the parent Hamiltonian.
\end{theorem}

\begin{proof}
    \leanok
    This follows immediately from Theorem~\ref{thm:friedrichs_gap_bound} by
    taking $\gamma$ to be the explicit constant produced there.
\end{proof}
